\documentclass[10pt,a4paper]{amsart}
\usepackage[margin=19mm]{geometry}
\usepackage{amsmath,amssymb,mathtools}
\usepackage[scr=rsfs,cal=euler]{mathalfa}
\usepackage{xfrac}
\usepackage[unicode,hidelinks,bookmarks=false]{hyperref}
\newcommand{\ie}{i.e.\ }
\newcommand{\cf}{cf.\ }
\newcommand{\resp}{respectively\ }
\newcommand{\Subsection}{\subsection}
\providecommand{\nobreakdash}{\nobreak\hbox{-}}
\setlength{\emergencystretch}{2em}
\multlinegap0pt
\usepackage{bm}
\usepackage{bbm}
\usepackage{dsfont}
\usepackage{xspace}
\usepackage{cancel}
\usepackage{mathtools}
\usepackage{amsthm}
\makeatletter
\@ifundefined{lemma}{\newtheorem{lemma}{Lemma}}{}
\makeatother
\title[Acyclic Poset Multiplihedra and their Quotients]{Acyclic Poset Multiplihedra and their Quotients}
\author{Stefan Forcey, Ross Glew, Christopher Tapo}
\date{Source: arXiv:2607.18651v2 (CC BY 4.0)}
\begin{document}
\maketitle

\noindent\emph{Source and licence.} Acyclic Poset Multiplihedra and their Quotients. Version arXiv:2607.18651v2 (CC BY 4.0), distributed under the Creative Commons Attribution 4.0 International licence, as recorded in the arXiv licence metadata for this exact version and shown on the abstract page https://arxiv.org/abs/2607.18651v2. Full rights notice: licenses/arxiv-2607.18651-CC-BY-4.0.txt.\par\smallskip

\noindent\textit{Editorial scope.} Counted unit: the Lemma whose source label is \texttt{None} in the article (source0.tex lines 486--488, source label \texttt{None}), delivered together with the article's own complete original proof of it (source0.tex lines 490--505, one \texttt{proof} environment of 16 lines). No mathematics is written, repaired or re-derived: statement and proof are the article's own text line for line. Required context retained uncounted: none. Every definition, display and label of those lines is retained unchanged. Omitted from this package and not redistributed: 1--485, 489--489, 506--788. Nothing inside a retained excerpt is omitted or altered: every retained body is delivered exactly as the article's lines are written, so no omission receipt of any kind accompanies this package. Citations: the retained excerpts carry no citation command at all, so no bibliography entry of the article is transcribed and no reference is redistributed. Reference adaptation: none was needed and none was made; every reference inside the delivered unit resolves inside it, so no unresolved reference or citation is produced. Printed slips kept literal: no printed word, symbol, hypothesis or number of the counted statement or of its proof was altered. Layout and class: the article's own class is \texttt{article} with packages including fontenc, hyperref, amsmath, amssymb, amsthm, orcidlink, booktabs, pdflscape, geometry, CJKutf8, microtype, babel, newpxtext, newpxmath; the wrapper is the lane's pinned \texttt{amsart} editorial wrapper at 10pt on A4, which restates the theorem-like environments the retained text needs and adds the article's own macro definitions that the retained text needs (none), with extra wrapper packages (bm, bbm, dsfont, xspace, cancel, mathtools). The delivered numbering is the unit's own. Assets: no retained span contains a figure environment, a graphics-inclusion command, a PDF-page inclusion, a table or a TikZ picture; no image file is copied into this package, the package declares no asset dependency and no image file is redistributed with it. Licence: version arXiv:2607.18651v2 of this article is distributed under the Creative Commons Attribution 4.0 International licence, as recorded in the arXiv licence metadata for this exact version and shown on the abstract page https://arxiv.org/abs/2607.18651v2. A full rights notice is delivered at licenses/arxiv-2607.18651-CC-BY-4.0.txt. Characters: every retained excerpt is pure ASCII, so the delivered engine, XeLaTeX with its default Latin Modern text font, reports no missing character.
 \begin{lemma}
     Let $F_T$ be a face of the graph multiplihedron $\mathcal{J}(L(P))$ associated to an acyclic tubing $T$ of the line graph of the Hasse diagram of $P$. Then $F_T$ intersects the evaluation space, which is the intersection of all the cycle hyperplanes.
 \end{lemma}       

   \begin{proof}
       %Let $F_T$ be a face of the multiplihedron  associated to an acyclic tubing $T$ of the line graph of the Hasse diagram of $P$.  The face $F_T$ is given by the intersection of the facet hyperplanes of the multiplihedron which contain that face, one for each tube of $T.$
       We consider an arrangement of cycle hyperplanes which determines the evaluation space. This arrangement is not unique, but we can choose one of the smallest such arrangements that determines the unique evaluation space.  
       The number of cycle hyperplanes is the same as the number of independent cycles in the (underlying graph of the) Hasse diagram, the cycle rank: $r = |E| - n +1$ for a connected diagram of $n$ nodes. (Kirchoff's cyclomatic number.) Notice that combining the equations of cycles gives a new cycle equation. We choose independent cycles by selecting a minimum set of edges in the Hasse diagram that need to be removed to make the underlying graph a tree, and then match a unique cycle to each of those edges, so that each cycle is broken by its matched edge. This choice provides a cycle basis, and the common intersection of the hyperplanes is the evaluation space. %Each face of the graph multiplihedron is determined by one equation for each tube. (There may be many more facets which contain that face, but one for each tube is sufficient to determine that face.) Thus the total number of equations in the realization of a face is at most $|E|,$ the number of edges of the Hasse diagram. 

       %Note: the equations themselves are of course not unique, they can be replaced by combinations of themselves with others, that is, row operations. We often add together independent cycles or subtract tube equations. Via the latter, we can get a set of equations for tubes where each variable has a non-zero coefficient only once, as in Definition 18 of \cite{dev-forc}.

 
Thus the cycle equalities we have chosen are linearly independent and there are fewer of them than the dimension we are working in, which is |E|, the number of edges in the Hasse diagram. Therefore the number of chambers is exactly $2^r$,  where $r =|E|-n+1$.  These chambers have all possible sign vectors of $r$ components from $\{+,-\}$. 
The number of red and blue pipes in $T$ is at most $n-1$, including the universal pipe, since we are considering an acyclic tubing of an $n$-node poset. To find any vertex of our face $F_T$ in the graph multiplihedron, combinatorially, we have to add tubes to $T$ until we get the max number of red and blue tubes, which is $|E|$. 
So in the case where we have a maximal acyclic poset tubing $T$, we need exactly $r$ more red and blue tubes to get $|E|$ of them. 

We claim that the final calculation of the coordinates from the maximal tubings will achieve all the combinations of $>$ and $<$ for the cycle equalities, reaching all $2^r$ chambers. To demonstrate our claim we need to add  tubes so that each choice of tube determines a $+$, or  $-$ , that is, the point obeys an inequality $>$ or $<$  for each of the cycle hyperplanes.   In the case where we have fewer than $n-1$ tubes in $T$, we start by arbitrarily adding more (and resolve some purple tubes) to get a maximal acyclic tubing.

 Given the maximal acyclic tubing $T$,  for each cycle $c$ of the Hasse diagram there will be a pair of edges of the Hasse diagram neither of which are the only edge of a single tube, but are either both in or both out of every tube in $T.$ Furthermore, in each pair, given a traversal of the cycle, one edge will be oriented with the direction of the traversal, and the other against. These pairs might overlap, sharing an edge between several pairs, specifically when cycles overlap. Choosing exactly one edge from each pair and deleting it produces a spanning tree; and this choice determines a cycle basis for the evaluation space. Taking that collection of cycles and considering its hyperplanes, we claim there is a vertex of $F_T$ from the graph multiplihedron in each chamber of the central arrangement. The claim follows since the collection of edges in the pairs mentioned are nodes in the line graph that can be completely ordered in nested depth via the additional tubes (in groupings of these nodes that correspond to cycles sharing an edge.) Furthermore, for each pair of edges we can choose the tubes to make either edge in the pair at a greater nested depth, and that is a decision independent of the options chosen for the other pairs.  Then for any hyperplane of the chosen cycle basis, we can (independently of the other hyperplanes) force the vertex to lie on either side of it, due to the exponential growth of the coordinate values as nested depth decreases. Examples follow; note that the color of the added tubes again preserves the inequalities needed since the blue tubes cannot be outside of red, and the red tubes confer an extra factor of 3. 
   \end{proof}  

\end{document}
